\clearpage
\fancyhead[L]{\small SECOND COMET OF 1805.}
\hypertarget{native-275-lower}{}
\begin{center}
SECOND COMET OF 1805.

\rule{0.6\linewidth}{0.4pt}

Monatliche Correspondenz. Vol. XIV. Pp. 75–86. July 1806.

\rule{0.6\linewidth}{0.4pt}

\textsc{Gauss to von Zach.}
\end{center}
\begin{flushright}
Brunswick, 20 May 1806.
\end{flushright}

\par\indent --- --- --- It has been quite some time since you were kind enough to send me, in manuscript, the original Marseille observations of the two comets that appeared last year. As soon as I received them, I took up the calculation of the second comet again, so I could have sent you the results two months ago. The delay was not due to neglect. I had meant to extend my investigation, and it was precisely this work that other commitments delayed and interrupted. I will no longer hesitate to share the investigations I have made so far. Though they have fallen short of the goal I originally wished and hoped to reach, they will still give rise to some reflections that may be of interest. With these preliminaries out of the way, I will now report on my calculations in order.

\par\indent Mr.~\textsc{Thulis}’s observations are all very good in themselves, but the stellar positions he observed with his equatorial instrument are mostly very inaccurate, especially the right ascensions. For example, the one on 29 November is too small by 25$^{\mathrm{s}}$ in time, the one on 6 December by 30$^{\mathrm{s}}$, and the one on 8 December by 38$^{\mathrm{s}}$. Some of the declinations also have to be changed by several minutes. It is no wonder, then, that many stars were marked with a note saying they could not be found, although I was able to reduce by far the greater part. I have reduced all the observations afresh except those of 11 and 22 November (for which I
\SourceLeaf{276}%
could not find the stars anywhere) and the one of 5 December, for which the star appears in \textsc{Bode} (64 Piscium) but nowhere in the Hist.\ cél. I therefore have the following additions to make about identifying the stars:

\begin{enumerate}
\item[1)] The stars of 15 and 16 November are the same star. Hist.\ cél., p.\ 477 (I do not have the Conn.\ des tems an XI).
\item[2)] The star of 23 November appears in the Hist.\ cél., p.\ 20.
\item[3)] The star on 3 December is 113 Piscium in \textsc{Bode}, and appears in the Hist.\ cél., pp.\ 37 and 200.
\item[4)] On 7 December it is 374 Aquarii in \textsc{Bode}, 981 in \textsc{Mayer}.
\item[5)] On 8 December it is 359 Aquarii in \textsc{Bode}, the same star with which Dr.\ \textsc{Olbers} and I compared the comet that evening.
\end{enumerate}

{\rightskip=0pt plus 1em
\par\indent Once I had carefully reduced the observations, I used them to improve my parabolic elements. It turned out that the three positions on which the calculation was based could not be represented by the parabola quite as accurately as one might have wished. The differences therefore had to be distributed as well as possible among all three observations. I did this with the greatest care, using a method of my own that I also intend to include in my book. This gave the following new parabolic elements for the second comet. Seeberg time. The motion is direct.\par}



\begin{center}\begingroup
\fontsize{9}{11}\selectfont
\setlength{\tabcolsep}{2pt}\renewcommand{\arraystretch}{1.12}
\begin{tabular}{lr}
Passage through perihelion, 1805 Dec. 31 & $6^{\mathrm{u}}\,54^{\mathrm{m}}\,34^{\mathrm{s}}$ \\
Longitude of perihelion .\ .\ .\ .\ .\ . & $109^{\circ}\,21'\,50''$ \\
Longitude of the ascending node .\ . & $250\;33\;35$ \\
Inclination of the orbit .\ .\ .\ .\ .\ .\ . & $16^{\circ}\,30'\,32''$ \\
Logarithm of the least distance .\ . & 9.9503300 \\
Logarithm of mean daily motion . & 0.0346333 \\
\end{tabular}\endgroup\end{center}



\par\indent Here is the comparison of these elements with all the observations:



\begin{center}\begingroup
\fontsize{8}{10}\selectfont
\setlength{\tabcolsep}{2pt}\renewcommand{\arraystretch}{1.12}
\begin{tabular}{|c|c|c|c|c|c|l|}
\hline
& & \shortstack{Calculated\\right\\ascension} & \shortstack{Calculated\\declination} & \multicolumn{2}{c|}{Difference in} & \\
1805 & Mean time & & & \shortstack{Right asc.\\R---B} & \shortstack{Declination\\R---B} & Observer \\
\hline
Nov.\ 15 & $10\,28^{\mathrm{m}}\,13^{\mathrm{s}}$ & $14^{\circ}\,33'\,0''$ & $39^{\circ}\,18'\,41''$ N. & $-2'\,30''$ & $+0'\,26''$ & Mr.\ Thulis \\
16 & $7\;\;17\;\;54$ & $14\;\;12\;\;7$ & $38\;\;58\;\;53$ & $-1\;\;58$ & $+0\;\;46$ & --- --- \\
19 & $7\;\;11\;\;17$ & $12\;\;55\;\;12$ & $37\;\;34\;\;27$ & $-0\;\;4$ & $+1\;\;7$ & --- --- \\
20 & $7\;\;3\;\;15$ & $12\;\;28\;\;23$ & $36\;\;59\;\;41$ & $-0\;\;23$ & $-0\;\;32$ & --- --- \\
21 & $9\;\;36\;\;1$ & $11\;\;57\;\;5$ & $36\;\;16\;\;16$ & $+1\;\;13$ & $-0\;\;20$ & --- --- \\
23 & $7\;\;23\;\;15$ & $11\;\;1\;\;15$ & $34\;\;46\;\;43$ & $+1\;\;0$ & $+0\;\;40$ & --- --- \\
29 & $6\;\;39\;\;44$ & $7\;\;24\;\;39$ & $26\;\;47\;\;49$ & $+1\;\;32$ & $+0\;\;1$ & --- --- \\
Dec.\ 2 & $5\;\;34\;\;9$ & $4\;\;48\;\;43$ & $18\;\;58\;\;18$ & $+1\;\;35$ & $-0\;\;52$ & Dr.\ Olbers \\
2 & $9\;\;39\;\;0$ & $4\;\;36\;\;3$ & $18\;\;20\;\;18$ & $+1\;\;47$ & $-0\;\;3$ & Mr.\ Thulis \\
3 & $5\;\;48\;\;39$ & $3\;\;39\;\;51$ & $15\;\;4\;\;52$ & $+0\;\;45$ & $-2\;\;5$ & Dr.\ Olbers \\
3 & $6\;\;30\;\;37$ & $3\;\;36\;\;52$ & $14\;\;55\;\;11$ & $+1\;\;10$ & $-0\;\;26$ & Mr.\ Thulis \\
4 & $6\;\;31\;\;5$ & $2\;\;16\;\;36$ & $10\;\;6\;\;29$ N. & $+2\;\;12$ & $-1\;\;12$ & --- --- \\
6 & $7\;\;16\;\;16$ & $358\;\;39\;\;22$ & $3\;\;39\;\;32$ S. & $+0\;\;22$ & $+0\;\;48$ & --- --- \\
7 & $7\;\;35\;\;21$ & $356\;\;9\;\;21$ & $13\;\;4\;\;56$ & $0\;\;0$ & $+0\;\;42$ & --- --- \\
8 & $5\;\;27\;\;22$ & $353\;\;20\;\;30$ & $22\;\;56\;\;24$ & $-1\;\;10$ & $+2\;\;44$ & Dr.\ Olbers \\
8 & $6\;\;46\;\;23$ & $353\;\;8\;\;33$ & $23\;\;34\;\;40$ & $-0\;\;42$ & $+2\;\;18$ & --- --- \\
8 & $6\;\;55\;\;17$ & $353\;\;7\;\;44$ & $23\;\;35\;\;37$ & $+0\;\;4$ & $-0\;\;47$ & Dr.\ Gauss \\
8 & & $353\;\;6\;\;21$ & $23\;\;40\;\;59$ & $-0\;\;33$ & $-0\;\;15$ & Dr.\ Maskelyne \\
8 & $6\;\;59\;\;39$ & $353\;\;3\;\;54$ & $23\;\;47\;\;34$ & $+0\;\;9$ & $-0\;\;33$ & Mr.\ Thulis \\
8 & $7\;\;58\;\;3$ & $352\;\;57\;\;34$ & $24\;\;6\;\;39$ & $-0\;\;23$ & .\quad . & Dr.\ Gauss \\
9 & $6\;\;4\;\;9$ & $349\;\;2\;\;44$ & $35\;\;23\;\;39$ & $-0\;\;16$ & .\quad . & Mr.\ Thulis \\
\hline
\end{tabular}\endgroup\end{center}



\par\indent The comparison was carried out with full precision, taking due
\SourceLeaf{277}%
account of aberration, nutation and parallax. The fractions of a second recorded in the observations for greater accuracy were dropped only at the end of the calculation. [The times are local to the places of observation.]

\par\indent Although the differences themselves are not much larger than observational errors commonly can be for comets, they follow an unmistakable pattern that seems to attest to the quality of \textsc{Thulis}’s observations. Indeed, in my own observation on 8 December, the nucleus shone through so distinctly that the comet must have been observable with considerable precision. It is much to be regretted that the only known meridian observations are those made at Greenwich on 8 December (reported in the newspapers) and at Marseille on the 9\textsuperscript{th}. From this regular pattern one could foresee that an orbit determined without the parabolic hypothesis, whether it turned out to be an ellipse or a hyperbola, would at least agree considerably better with the observations. I therefore undertook the work all the more gladly because I was then working on some important changes and additions to my method and wanted a practical application for them. This was especially welcome because \emph{this} example has a peculiarity that cannot occur with the asteroids Ceres, Pallas and Juno. This is no place to go into it at greater length.

\par\indent The result of this work is indeed remarkable in more than one respect. I obtained the following \emph{elliptical} orbit:

\begin{center}\begingroup
\fontsize{9}{11}\selectfont
\setlength{\tabcolsep}{2pt}\renewcommand{\arraystretch}{1.12}
\begin{tabular}{lr}
\multicolumn{2}{l}{Passage through perihelion} \\
\quad on 2 January 1806, Seeberg time .\ . & $11^{\mathrm{u}}\;8^{\mathrm{m}}\;45^{\mathrm{s}}$ \\
Longitude of perihelion .\ .\ .\ .\ . & $109^\circ\;30'\;2''3$ \\
Longitude of the asc. node .\ . & $251\;28\;22.5$ \\
Inclination of the orbit .\ .\ .\ .\ .\ .\ . & $12\;43\;10.0$ \\
Logarithm of the semimajor axis .\ .\ . & $0.4505887$ \\
Mean daily motion .\ .\ .\ .\ . & $748''383$ \\
Period of revolution .\ .\ .\ .\ .\ .\ .\ . & 1731 days 17 hr. \\
Eccentricity .\ .\ .\ .\ .\ .\ .\ .\ .\ . & $0.6769242$ \\
Greatest distance from the Sun .\ .\ .\ . & $4.732625$ \\
Least \quad---\ ---\ --- .\ .\ .\ . & $0.911786$ \\
Logarithm of the least distance .\ .\ . & $9.9598931$ \\
\end{tabular}\endgroup\end{center}

\Needspace{28\baselineskip}
\par\indent These elements agree with the observations as follows:

\begin{center}\begingroup
\fontsize{9}{11}\selectfont
\setlength{\tabcolsep}{2pt}\renewcommand{\arraystretch}{1.12}
\begin{tabular}{c|cc|cc|l}
& \shortstack{Calculated right\\ascension} & \shortstack{Calculated\\declination} & \multicolumn{2}{c|}{Difference in} & \\
1805 & & & \shortstack{Right asc.\\R---B} & \shortstack{Declination\\R---B} & Observer \\
\hline
Nov. 15 & $14^\circ\;35'\;45''$ & $39^\circ\;17'\;56''$ N. & $+15''$ & $-19''$ & Mr. Thulis \\
16 & $14\;14\;20$ & $38\;56\;6$ & $+15$ & $-1$ & --- --- \\
19 & $12\;55\;23$ & $37\;34\;5$ & $+7$ & $+45$ & --- --- \\
20 & $12\;27\;59$ & $36\;59\;24$ & $-47$ & $-49$ & --- --- \\
21 & $11\;55\;59$ & $36\;15\;58$ & $+7$ & $-38$ & --- --- \\
23 & $10\;59\;33$ & $34\;46\;25$ & $-42$ & $+22$ & --- --- \\
29 & $7\;22\;34$ & $26\;47\;40$ & $-33$ & $-8$ & --- --- \\
Dec.\ 2 & $4\;47\;28$ & $18\;58\;28$ & $+20$ & $-42$ & Dr. Olbers \\
2 & $4\;34\;22$ & $18\;20\;44$ & $+6$ & $+13$ & Mr. Thulis \\
3 & $3\;38\;51$ & $15\;5\;11$ & $-15$ & $-1'\;46$ & Dr. Olbers \\
3 & $1\;35\;50$ & $14\;55\;37$ & $+8$ & $0$ & Mr. Thulis \\
4 & $2\;15\;55$ & $10\;7\;4$ N. & $+1'\;31$ & $-37$ & --- --- \\
6 & $358\;39\;13$ & $3\;38\;58$ S. & $+13$ & $+14$ & --- --- \\
7 & $356\;9\;28$ & $13\;4\;43$ & $+7$ & $+29$ & --- --- \\
8 & $353\;21\;10$ & $22\;56\;53$ & $-30$ & $+3\;13$ & Dr. Olbers \\
8 & $353\;9\;12$ & $23\;35\;19$ & $-3$ & $+2\;57$ & --- --- \\
8 & $353\;8\;10$ & $23\;36\;16$ & $+30$ & $-8$ & Dr. Gauss \\
8 & $353\;6\;55$ & $23\;41\;36$ & $+1$ & $+22$ & Dr. Maskelyne \\
8 & $353\;4\;18$ & $23\;48\;8$ & $+33$ & $+1$ & Mr. Thulis \\
8 & $352\;57\;46$ & $24\;7\;18$ & $-11$ & .\quad . & Dr. Gauss \\
9 & $349\;2\;56$ & $35\;25\;21$ & $-4$ & .\quad . & Mr. Thulis \\
\end{tabular}\endgroup\end{center}
\SourceLeaf{278}%
\par\indent Comparing these differences with those found above for the parabola, one sees at once that the ellipse fits \emph{very much} better. Among \textsc{Thulis}’s observations, the right ascension on 4 December is the only one for which the difference exceeds $50''$. The star used that day was one I took from the Conn. des tems, An. XIII, but I could find no observations of it in the Hist. cél. It may need a correction. Dr. \textsc{Olbers} himself had described his declination on 3 December as very doubtful. On 8 December, it was impossible to bring the calculation into agreement with Dr. \textsc{Olbers}’s declinations without departing by 3 minutes from the Greenwich meridian observation, with which \textsc{Thulis} and I also agree except for slight differences. The declination on 9 December, when reduced to its apparent value by applying refraction, also agrees better with \textsc{Thulis}’s estimate in the ellipse than in the parabola.

\par\indent The following conclusions can now be drawn:

\par\indent The observations of the comet can be represented by all the ellipses that lead from the one above to the parabola. More explicitly, an ellipse with \emph{any} mean distance greater than 2.82 still represents them better than any parabola, though the fit gradually worsens the farther it is from the ellipse above. On the other side too, down to a certain limit, a mean distance smaller than 2.82 will still represent the observations better than the parabola does, since the differences increase only gradually on either side of the minimum. I have made no calculation on this point, but I believe that even a mean distance of 2.0, or a period of revolution of less than three years, would still agree just as well with all the observations as the parabola does.

\par\indent I have often thought and said that, among the hundred comets calculated so far, for which the parabolic hypothesis is routinely applied without further inquiry and considered satisfactory if it agrees with the observations, probably by far the greater part are such that the observations alone would not even prove with certainty that their orbits must come very close to a parabola. For every comet whose orbit is calculated, it would therefore be desirable to have some idea how far apart the bounds are within which the orbit must lie if the calculation is to stay reasonably close to the observations. In this respect our comet, which in 24 days showed a geocentric motion of nearly 80 degrees, is in my opinion a very striking example of how uncertain we remain about the true dimensions of the orbit. Still, I readily grant that if the comet had continued to be well observed into January (which could quite readily have been done in the southern hemisphere), it should have been possible to say something fairly definite about its true orbit. It would already be welcome in this regard if you could obtain \textsc{Thulis}’s original observations of 10 and 12 November for me, and sooner or later a precise determination of the stars used for comparison on 11 November. Extending the period of observation by five days is already of some value here.

\par\indent Despite all this uncertainty, the much better fit of the ellipse above did seem to make it somewhat probable that this comet really moves in an ellipse with a period of revolution that is not too long. There was also the resemblance to the comet of 1772. \textsc{Bessel} had found by calculation that if each of these comets was represented by an ellipse with a period of 33 years, the other elements showed no greater agreement than with the parabolas. Nevertheless, after the investigation above, it was still possible that one and the same comet might have completed, in those 33 years, two, three
\SourceLeaf{279}%
or more revolutions. This faint comet had been so conspicuous last year only because it came so close. Its faintness would readily explain why it went unnoticed on the other passages through perihelion, which may have occurred under unfavorable circumstances. My calculations on this point show, however, that this will not work either. I find that the same least and greatest distances from the Sun cannot be fitted to the observations of both comets at once without producing too large a difference in the other elements, such as the longitude of perihelion, the inclination of the orbit and the longitude of the node.

\par\indent Here is the table of two sets of elements for the comet of 1772. The first is for a parabola in which I took the least distance to be the same as in my (earlier) parabolic elements of last year’s comet. The second is for an ellipse with the same least and greatest distances as the ellipse above. Comparing the corresponding elements with those of last year’s comet shows that the comet of 1772, if its greatest and least distances are taken to be the same as for the comet of 1805, always has a considerably smaller longitude of perihelion, a considerably greater longitude of the node and an inclination greater by $4^\circ$ to $5^\circ$ than the later comet.

\begin{center}\begingroup
\fontsize{8}{10}\selectfont
\setlength{\tabcolsep}{2pt}\renewcommand{\arraystretch}{1.12}
\begin{tabular}{l|c|c}
& \multicolumn{2}{c}{Elements of the comet of 1772.} \\
 & Parabolic & Elliptical \\
\hline
Passage through perihelion & 9 Feb. 1772. $5^{\mathrm{u}}$ & 8 Feb. 1772. $1^{\mathrm{u}}$ \\
Longitude of perihelion . . . . & $90^{\circ}\,17'$ & $97^{\circ}\,21'$ \\
Longitude of the ascending node & $261\ \ 9$ & $263\ \ 24$ \\
Inclination of the orbit . . . . . & $20\ \ 28$ & $17\ \ 39$ \\
Least distance . . . . . . & $0.8918$ & $0.9118$ \\
Greatest distance . . . . . . & $\infty$ & $4.7326$ \\
\end{tabular}\endgroup\end{center}

\par\indent These elements are based only on the first and last of \textsc{Messier}’s observations (two observations always suffice if two of the elements are taken as given). The middle observation then gave a difference in longitude of 6 minutes for the parabola and 3 min. for the ellipse. The latitude agreed well in both cases. By distributing the differences, however, one could reduce them to about half without \emph{greatly} changing the elements themselves.

\par\indent Even so, I am not yet willing to rule out the possibility that the two comets are identical. The investigations so far do show that their motions cannot both be represented satisfactorily by a single orbit. There remains, however, the possibility that an external perturbation changed the elements considerably. If the comet of 1772 is assigned a period of 33 years, so that it had to return to perihelion in 1805, no close approach to a major planet is found. But the ellipse from whose perihelion it departed in 1772 and the one at whose perihelion it arrived in 1805 may still have quite different periods of revolution. In the intervening time, it may have passed so close to a major planet that the planet could have transformed the first ellipse into the second, whether the comet completed one revolution or several during that time. There remains, then, the admittedly rather lengthy investigation of whether, among all the elliptical orbits that can represent the motions of the two comets (and that, as we know for certain, are confined within such wide bounds), there are two such that each orbit gives its comet the very same position on some day between 1772 and 1805, and a major planet was very close to that position on that very day. In such a case the major planet would have had a strong
\SourceLeaf{280}%
influence on the elements, but its effect would have varied considerably with the relative motion. A change in the elements of whatever kind could thus be readily explained. I intend to investigate this further in the future. --- --- ---